\documentclass[10pt,a4paper]{amsart}
\usepackage[margin=19mm]{geometry}
\usepackage{amsmath,amssymb,mathtools}
\usepackage[scr=rsfs,cal=euler]{mathalfa}
\usepackage{xfrac}
\usepackage[unicode,hidelinks,bookmarks=false]{hyperref}
\newcommand{\ie}{i.e.\ }
\newcommand{\cf}{cf.\ }
\newcommand{\resp}{respectively\ }
\newcommand{\Subsection}{\subsection}
\providecommand{\nobreakdash}{\nobreak\hbox{-}}
\setlength{\emergencystretch}{2em}
\multlinegap0pt
\usepackage{bm}
\usepackage[normalem]{ulem}
\newtheorem{lemma}{Lemma}
\newtheorem{theorem}{Theorem}
\newcommand{\assign}{\mathrel{\mathop:}=}
\newcommand{\scS}{\mathcal{S}}
\newcommand{\txtb}{\textup{\texttt{b}}}
\newcommand{\txtc}{\textup{\texttt{c}}}
\newcommand{\wV}{\widetilde{V}}
\newcommand{\whV}{\widehat{V}}
\newcommand{\whp}{\widehat{p}}
\newcommand{\whv}{\widehat{v}}
\newcommand{\whw}{\widehat{w}}
\title{Spectral coarse spaces based on indefinite operators: the $H_k$-GenEO method}
\author{Th\'eophile Chaumont-Frelet, Victorita Dolean, Mark Fry, Ivan G. Graham, Matthias Langer}
\date{Source: arXiv:2605.31552v1 (CC BY 4.0)}
\begin{document}
\maketitle

\noindent\emph{Source and licence.} Spectral coarse spaces based on indefinite operators: the $H_k$-GenEO method. Version arXiv:2605.31552v1 (CC BY 4.0), distributed by arXiv under the Creative Commons Attribution 4.0 International licence, http://creativecommons.org/licenses/by/4.0/, as recorded in the arXiv licence metadata for this exact version and shown on the abstract page https://arxiv.org/abs/2605.31552v1. Full licence notice: \texttt{licenses/arxiv-2605.31552-CC-BY-4.0.txt}.\par\smallskip

\noindent\textit{Editorial scope.} Counted unit: the article's theorem lem\_ind\_ev (source0.tex lines 746--772). The article's complete original proof of it is delivered with it (source0.tex lines 774--852, one \texttt{proof} environment of 79 lines). No mathematics is written, repaired or re-derived: the statement and the proof are the article's own lines, in the article's own notation, and no retained line is substituted or re-flowed.  Retained uncounted as the source-local context the proof needs: lines 731--738.  Nothing was omitted from any retained span, so the package carries no omission record.  No retained line contains a citation, so the wrapper carries no bibliography and no bibliography entry.  No retained line contains a figure environment, an image-inclusion command or drawing code, so no image is delivered and the package declares no asset dependency.  Editorial formatting only, in addition: an AMS \texttt{amsart} preamble that restates the article's own declarations and macros used by the retained lines, namely the environments \texttt{lemma}, \texttt{theorem} on separate counters, together with the standard packages those lines need.  The retained environments print in this document as Lemma 1, Theorem 1 in order of appearance, whereas the article prints the same items with the numbering of its own full document.  The complete native source of the article as distributed in the arXiv e-print for this exact version is bundled unchanged as \texttt{sources/201730/source0.tex}.
\begin{lemma} \label{cpd}
Suppose $\wV$ is a real vector space with $\dim \wV = \tilde n$.   
Let $\txtc$ be a symmetric positive semi-definite bilinear form on $\wV$, 
let $\ker \txtc \assign \{ v \in \wV: \txtc(v,w) = 0 \ \text{for all} \ w \in \wV\}$ 
and suppose that $\dim \ker \txtc = s < \tilde n$. 
Further, let $\whV\subseteq \wV$ be an $(\tilde n-s)$-dimensional subspace such that $\whV\cap\ker\txtc=\{0\}$.
Then $\wV=\whV\oplus\ker\txtc$, where $\oplus$ denotes a direct sum, and $\txtc$ is positive definite on $\whV$.
\end{lemma}

\begin{theorem}[Abstract theorem for indefinite eigenvalue problems]
\label{lem_ind_ev}
Let $\wV$, $\txtc$ and $s$ be 
as in Lemma~\ref{cpd} and suppose also that $\txtb$ is a symmetric \textup{(}possibly indefinite\textup{)} bilinear form on $\wV$. 
Consider the generalised eigenvalue problem: find $\lambda \in \mathbb{R}$ and $p \in \wV \backslash \{0\}$ such that 
\begin{equation} \label{Ivan1}
  \txtb(p,v)  = \lambda \txtc(p,v) \qquad \text{for all} \quad v \in \wV.
\end{equation}
If $s \ge 1$, suppose also that $\txtb$ is positive definite on $\ker\txtc$. 
Then  $\wV$ has a basis $\{p_j: j = 1, \ldots,\tilde n\}$,  in which  $\{p_j: j = 1, \ldots, \tilde n-s\}$ are eigenvectors of \eqref{Ivan1}  which
can be chosen to be orthonormal with respect to $\txtc$ and the corresponding finite eigenvalues can be ordered
\begin{equation} \label{Ivan11}
	-\infty < \lambda_1 \leq  \lambda_2 \leq   \ldots \leq \lambda_{\tilde n-s} < \infty. 
\end{equation} 
The remaining  basis vectors $\{p_j: j = \tilde n-s+1, \ldots, \tilde n\}$  form a basis of $\ker\txtc$. 
In this sense we say that they correspond to infinite  eigenvalues of \eqref{Ivan1}.
Moreover, 
\begin{alignat}{2}
	\txtb(p_j,p_{j'}) &= \lambda_j \delta_{j,j'}, \quad && \ j \in \{1, \ldots, \tilde n-s\}, \ j'\in\{1,\ldots,\tilde n\};
	\label{orth_b}
	\\
	\txtc(p_j,p_{j'}) &= \delta_{j,j'}, \quad && \ j\in\{1,\ldots,\tilde n-s\},\ j'\in\{1,\ldots,\tilde n\}.
	\label{orth_c}
\end{alignat}
If $s=0$ (trivial case) the same conclusion holds but the positive definiteness condition on $\txtb$ is not required, 
and all eigenvalues are finite.
\end{theorem}

\begin{proof}
Suppose $s \ge 1$ and $\txtb$ is positive definite on $\ker \txtc$. 
Let us choose $\whV\subseteq\wV$ such that $\wV = \whV\oplus \ker \txtc$ (as in Lemma~\ref{cpd});
then every $p \in \wV$ has a unique decomposition $p = \whp + p_{\txtc} $ with $\whp\in \whV$ and  $p_{\txtc}\in \ker \txtc$. 
Thus the GEVP \eqref{Ivan1} is equivalent to seeking $\lambda \in \mathbb{R}$, $\whp \in \whV$ and  $p_{\txtc} \in \ker \txtc$ such
that $\whp + p_{\txtc} \ne 0$ and  
\begin{align} \label{Ivan1a} 
	\txtb(\whp + p_{\txtc}, v ) = \lambda \txtc (\whp, v) \qquad \text{for all} \;\; v \in \wV.
\end{align}
This, in turn, is equivalent to solving  the coupled system
\begin{alignat}{2}
	\txtb(\whp, v_\txtc) + \txtb(p_{\txtc}, v_{\txtc})    &= 0 \qquad 
	&& \text{for all} \;\; v_{\txtc} \in \ker \txtc, 
	\label{Ivan1b}\\
	\txtb(\whp,\whv) + \txtb(p_{\txtc},\whv) &= \lambda \txtc (\whp,\whv) \qquad\quad 
	&& \text{for all} \;\; \whv \in \whV,
	\label{Ivan1c}
\end{alignat}
for $\lambda \in \mathbb{R}$,   $\whp\in \whV$ and $p_{\txtc} \in \ker \txtc$ such
that $ \whp + p_{\txtc} \ne 0$. 
    
Now, by assumption on  $\txtb$, for each $\whw \in \whV$, there exists a unique $w_\txtc \in \ker \txtc$ such that
\begin{align}\label{Ivan1c1} 
	\txtb(w_\txtc, v_\txtc) = - \txtb(\whw , v_\txtc) \qquad \text{for all} \;\; v_\txtc \in \ker \txtc. 
\end{align} 
Denoting the solution to \eqref{Ivan1c1} as $w_\txtc = \scS \whw$, this defines a  linear map
$\scS: \whV\rightarrow \ker \txtc$ with the property
\begin{align}\label{Ivan1b1}
	\txtb(\scS \whw, v_\txtc) = - \txtb(\whw, v_\txtc) \qquad \text{for all} \;\;
	\whw \in \whV \;\; \text{and} \;\; v_\txtc \in \ker \txtc,
\end{align}
or, equivalently,  
\begin{align} \label{Ivan1d} 
	\txtb\bigl((I + \scS) \whw, v_\txtc\bigr) = 0 \qquad \text{for all} \;\; \whw \in \whV \;\; \text{and} \;\; v_\txtc \in \ker \txtc.
\end{align} 

We now  consider the GEVP: find $\lambda \in \mathbb{C}$ and   $\whp \in \whV\backslash \{0\}$ such that
\begin{align} \label{Ivan1e} 
	\widetilde{\txtb} (\whp, \whv) \assign \txtb\bigl((I + \scS) \whp, \whv\bigr) = \lambda \txtc(\whp, \whv) \qquad 
	\text{for all} \;\; \whv \in \whV.  
\end{align}
(This is the Schur complement of \eqref{Ivan1b}, \eqref{Ivan1c}.)
    
Note that $\scS$ is self-adjoint with respect to $\txtb$. This is because, using the symmetry of $\txtb$ (at the first and third steps) 
and the definition of $\scS$ (at the second and fourth steps),  
\[
	\txtb(\scS \whp, \whv) = \txtb(\whv, \scS \whp) = -\txtb(\scS \whv, \scS \whp) 
	= -\txtb( \scS \whp,\scS \whv) =  \txtb(\whp, \scS \whv).
\] 
Thus the bilinear form $\widetilde{\txtb}$ on the left-hand side of \eqref{Ivan1e} is symmetric on $\whV$.
Now $\txtc$ is symmetric and (by Lemma~\ref{cpd}) positive definite on $\whV$;
so the GEVP \eqref{Ivan1e} has $\tilde n-s$ eigenpairs,
which we denote  $(\lambda_j, \whp_{j})$, $j = 1, \ldots , \tilde n-s$, with real finite eigenvalues $\lambda_j$,
which can be ordered as in \eqref{Ivan11}.  The $\whp_{j} \in \whV$ can be chosen orthonormal with respect to $\txtc$ 
and form a basis of $\whV$.  
Now we define 
\begin{align} \label{eigdefj} 
	p_j = (I + \scS)\whp_{j}, \qquad j = 1, \ldots, \tilde n-s.
\end{align} 
By \eqref{eigdefj}, \eqref{Ivan1e} and since $\scS \whp_{j} \in \ker \txtc$, we have 
\[
	\txtb(p_j, \whv) = {\widetilde{\txtb}(\whp_{j},\whv)}
	= \lambda_j \txtc(\whp_{j}, \whv) = \lambda_j \txtc(p_{j}, \whv)
	\qquad \text{for all} \;\; \whv \in \whV.
\]
Also, by \eqref{eigdefj} and \eqref{Ivan1d} we have $\txtb(p_j, v_{\txtc}) = 0$ for all $v_{\txtc} \in \ker \txtc$, and so
(see \eqref{Ivan1b} and  \eqref{Ivan1c}), 
$(\lambda_j, p_j),\,  j = 1, \ldots ,\tilde n-s$
are eigenpairs of the original GEVP \eqref{Ivan1}.
Now, choosing an arbitrary basis, $\{p_{\tilde n-s+1},\ldots,p_{\tilde n}\}$, of {$\ker\txtc$} we obtain a basis of $\wV$. 
For $j \in \{1, \ldots , \tilde n-s\}$, relation \eqref{orth_c} follows for $ j' \in \{ 1, \ldots, \tilde n-s\}$ 
because  $\scS \whp_{j}, \scS \whp_{j'} \in \ker \txtc$, and so $\txtc(p_j,p_{j'}) = \txtc(\whp_{j},\whp_{j'}) = \delta_{j,j'}$ 
since the $\{\whp_{j}: j = 1, \ldots, \tilde n-s\}$ were chosen orthonormal with respect to $\txtc$.   
For $j' \in  \{\tilde n-s+1, \ldots, \tilde n\}$,  \eqref{orth_c} is trivial because then $p_{j'} \in \ker \txtc$.

Relation \eqref{orth_b} then follows directly  from \eqref{orth_c} because,  for $j \in \{1, \ldots, \tilde n-s\}$ 
and any $j' {\in \{1,\dots,\tilde n\}}$, $\txtb(p_j, p_{j'}) = \lambda_j\txtc(p_j, p_{j'})$. 
This completes the proof when $s\ge1$, while the case $s = 0$ is trivial since \eqref{Ivan1} can then be reduced to a standard eigenvalue problem.     
\end{proof}

\end{document}
